\section{Results}
\label{sec:results}

\subsection{End-to-End Document Quality}

Table~\ref{tab:quality} summarises the mean expected document quality $\bar{Q}(s)$
at each pipeline stage completion event for the three configurations: Full Pipeline,
No-Refs, and No-Visuals.
Standard errors are reported at the 95\% confidence level.

\begin{table}[h!]
\centering
\caption{Mean document quality $\bar{Q}(s)$ at each pipeline stage ($n=120$ sessions).
         95\% confidence intervals in parentheses.}
\label{tab:quality}
\begin{tabular}{lrrr}
\toprule
Stage & Full Pipeline & No-Refs & No-Visuals \\
\midrule
$s_0$ (Initial)        & 0.05 (0.05, 0.05) & 0.05 & 0.05 \\
$s_1$ (Plan)           & 0.18 (0.16, 0.20) & 0.15 & 0.18 \\
$s_2$ (Extr.+Verify)   & 0.45 (0.43, 0.47) & 0.20 & 0.42 \\
$s_3$ (Draft)          & 0.78 (0.76, 0.80) & 0.55 & 0.70 \\
$s_4$ (Assemble)       & 0.91 (0.89, 0.93) & 0.68 & 0.82 \\
$s_5$ (Validate)       & \textbf{0.97} (0.96, 0.98) & 0.75 & 0.88 \\
\bottomrule
\end{tabular}
\end{table}

The full pipeline achieves $\bar{Q}(s_5)=0.97$, a relative improvement of
$29.3\%$ over No-Refs ($0.75$) and $10.2\%$ over No-Visuals ($0.88$).
The quality S-curve is visualised in Figure~\ref{fig:quality_scurve}, where the
shaded region between the Full Pipeline and No-Refs curves quantifies the quality
contribution of the reference verification subsystem at each stage.

\begin{figure}[h!]
\centering
\begin{tikzpicture}
\begin{axis}[PL,
  xlabel={Pipeline stages completed},
  ylabel={Expected document quality $Q(s)$},
  title={Cumulative quality gain per pipeline stage},
  xmin=0, xmax=5.5, ymin=0, ymax=1.05,
  xtick={0,1,2,3,4,5},
  xticklabels={$s_0$,Plan,Extr.+\\Verify,Draft,Assemble,Validate},
  x tick label style={align=center, font=\tiny},
  ytick={0,0.2,0.4,0.6,0.8,1.0},
  legend pos=north west,
  width=13cm, height=7cm,
]
\addplot[name path=FULL, C1, very thick, mark=*, mark options={fill=C1,solid}]
  coordinates{(0,0.05)(1,0.18)(2,0.45)(3,0.78)(4,0.91)(5,0.97)};
\addlegendentry{Full pipeline (refs + visuals)}
\addplot[name path=NOREF, C4, thick, dashed, mark=square*, mark options={fill=C4,solid}]
  coordinates{(0,0.05)(1,0.15)(2,0.20)(3,0.55)(4,0.68)(5,0.75)};
\addlegendentry{No references}
\addplot[name path=NOVIS, C3, thick, dashed, mark=triangle*, mark options={fill=C3,solid}]
  coordinates{(0,0.05)(1,0.18)(2,0.42)(3,0.70)(4,0.82)(5,0.88)};
\addlegendentry{No visuals}
\addplot[F1, opacity=0.40] fill between[of=FULL and NOREF];
\foreach \s in {1,2,3,4,5}{\addplot[DG!20,thin]coordinates{(\s,0)(\s,1.03)};}
\end{axis}
\end{tikzpicture}
\caption{Cumulative document quality as pipeline stages complete. Shaded area
         quantifies the quality gain attributable to the reference verification
         subsystem. The drafting stage ($s_2\to s_3$) contributes the largest
         single quality increment across all configurations.}
\label{fig:quality_scurve}
\end{figure}


\subsubsection*{Stage-by-Stage Gain Decomposition}

The largest single-stage quality increment in the full pipeline is at $s_2\to s_3$
(the Draft stage), with $\Delta Q_3 = 0.78 - 0.45 = 0.33$.
This gain is larger than any other single-stage increment and is consistent across
all five topic categories in the corpus (range: $0.29$--$0.37$).
The gain at $s_2\to s_3$ is primarily attributable to the $R(s)$ sub-metric:
after $F_2$ admits only \texttt{enhanced} references, the drafting stage generates
prose that is factually grounded in verified sources, causing $R(s_3)$ to jump
from the post-verification floor to a near-ceiling value.

The second-largest increment is at $s_1\to s_2$ (Extract and Verify), with
$\Delta Q_2 = 0.45 - 0.18 = 0.27$.
This is entirely attributable to $R(s_2)$: as the \cgrag{} gate processes candidate
references and builds $\mathcal{R}^+$, $R(s)$ rises monotonically.
The No-Refs ablation shows $\Delta Q_2 = 0.20 - 0.15 = 0.05$ over the same interval,
confirming that almost all of the $s_1\to s_2$ gain in the full pipeline comes from
reference verification rather than from outline refinement.

The assembly stage $s_3\to s_4$ contributes $\Delta Q_4 = 0.91 - 0.78 = 0.13$,
driven primarily by the $V(s)$ sub-metric: as figures are embedded in the document
and cross-references are resolved, visual coherence improves.
The No-Visuals ablation shows $\Delta Q_4 = 0.82 - 0.70 = 0.12$, which is similar
in magnitude, indicating that the assembly stage's $V(s)$ gain is comparable
between configurations because even the No-Visuals document benefits from
the resolution of section cross-references (as opposed to figure cross-references).

The validation stage $s_4\to s_5$ contributes $\Delta Q_5 = 0.97 - 0.91 = 0.06$,
driven by compilability improvements: the \tvg{} repair cycle at validation time
fixes any residual TikZ errors, increasing $C(s_5)$ from $0.95$ (pre-validation)
to $0.99$ (post-validation).

\subsubsection*{Topic-Category Breakdown}

Table~\ref{tab:quality_by_topic} reports $\bar{Q}(s_5)$ broken down by topic category.

\begin{table}[h!]
\centering
\caption{Mean document quality $\bar{Q}(s_5)$ by topic category (full pipeline only).}
\label{tab:quality_by_topic}
\begin{tabular}{lrrl}
\toprule
Topic Category & $n$ & $\bar{Q}(s_5)$ & Notes \\
\midrule
CS \& Software Engineering & 38 & 0.98 & Highest $C(s_5)$, rich TikZ support \\
Education Research         & 26 & 0.97 & Average across all sub-metrics \\
Biomedical Informatics     & 21 & 0.96 & Lower $R(s_5)$: preprint-heavy domain \\
Social Sciences            & 19 & 0.97 & Slightly lower $V(s_5)$ (fewer figures) \\
Engineering                & 16 & 0.96 & More complex figures, lower $C(s_5)$ \\
\midrule
Overall                    & 120 & 0.97 & \\
\bottomrule
\end{tabular}
\end{table}

The variation across topic categories is small (range: $0.96$--$0.98$), indicating
that the formal system design provides consistent quality regardless of domain.
The two categories with $\bar{Q}(s_5)=0.96$ (biomedical and engineering) have
identifiable sub-metric weaknesses: biomedical sessions have a higher fraction
of preprint-level ($Q_3$) references because much recent biomedical research
is published first on bioRxiv, which lowers the median $\mathcal{L}_Q$ tier
and therefore $R(s_5)$; engineering sessions produce more complex TikZ diagrams
(multi-layer schematics, detailed flow charts) that have a slightly lower first-attempt
success probability $p_{\mathrm{eng}}=0.67$ versus the corpus-wide $\hat{p}=0.72$.

\subsection{Invariant Verification}

Table~\ref{tab:invariants} reports the empirical invariant satisfaction rates across
all 120 sessions.

\begin{table}[h!]
\centering
\caption{System invariant satisfaction rates ($n=120$ sessions).
         ``By construction'' means the invariant holds provably, not empirically.}
\label{tab:invariants}
\begin{tabular}{llll}
\toprule
Invariant & Description & Rate & Status \\
\midrule
I1 & Preamble Uniqueness   & 120/120 (100\%) & By construction \\
I2 & Citation Safety       & 120/120 (100\%) & By construction \\
I3 & Billing Monotonicity  & 120/120 (100\%) & By construction \\
I4 & Order Preservation    & 120/120 (100\%) & By construction \\
I5 & Real-Data Guarantee   & 118/120 (98.3\%) & Infrastructure-conditional \\
\bottomrule
\end{tabular}
\end{table}

Invariants I1--I4 hold with 100\% satisfaction across all 120 sessions.
Their status as ``by construction'' reflects that their correctness follows from the
proofs given in Section~\ref{sec:architecture}, not from empirical verification:
we would not expect to find violations unless a code change had broken the formal
invariant, in which case the violation count would be bounded by the number of
sessions run after the breaking change.
No such changes occurred during the evaluation period.

Invariant~I5 was violated in 2 of 120 sessions (98.3\% satisfaction rate).
Investigation revealed that both violations occurred on the same day (evaluation day 23),
during a 14-minute window in which the Semantic Scholar API endpoint returned stale
responses served from an intermediary CDN cache.
The \texttt{If-Modified-Since} header was present in the request, but the CDN
did not respect it and returned cached responses without updating the
\texttt{Last-Modified} header.
The references admitted in these two sessions are therefore not guaranteed to have
been live-fetched at session time $t_0$.

We classify these as infrastructure failures (CDN misconfiguration at the upstream
provider) rather than design failures.
The formal specification of I5 (Section~\ref{sec:architecture}) is conditional
on network liveness; the CDN failure caused a violation of that condition.
A future version of the verification protocol will add a content-hash check
against a freshness oracle to detect CDN caching anomalies.

\subsection{Repair Convergence}

Table~\ref{tab:repair} presents the empirical repair convergence data alongside
the theoretical predictions from Corollary~\ref{cor:repair}.

\begin{table}[h!]
\centering
\caption{Cumulative TikZ figure compilation success by attempt number.
         Theoretical values from Corollary~\ref{cor:repair} with $\hat{p}=0.72$.
         $n_{\mathrm{total}}=480$ figures.}
\label{tab:repair}
\begin{tabular}{lrrrrrr}
\toprule
\multirow{2}{*}{Attempt} &
\multicolumn{2}{c}{Empirical} &
\multicolumn{2}{c}{Theoretical} &
\multirow{2}{*}{Diff (pp)} \\
\cmidrule(lr){2-3}\cmidrule(lr){4-5}
& Count & \% & $P(\cdot\mid K)$ & \% & \\
\midrule
$K=0$ (1st attempt)  & 335 & 69.8 & $1-0.28^1=0.72$ & 72.0 & $-2.2$ \\
$K=1$ (2nd attempt)  & 427 & 89.0 & $1-0.28^2=0.9216$ & 92.2 & $-3.2$ \\
$K=2$ (3rd attempt)  & 470 & 97.9 & $1-0.28^3=0.9780$ & 97.8 & $+0.1$ \\
$K=3$ (4th attempt)  & 477 & 99.4 & $1-0.28^4=0.9938$ & 99.4 & $0.0$ \\
\bottomrule
\end{tabular}
\end{table}

Theoretical predictions from Corollary~\ref{cor:repair} agree with empirical results
to within 3.2 percentage points at every attempt level.
The small systematic negative bias at $K=0$ and $K=1$ (empirical slightly below
theoretical) is attributable to a cluster of 8 figures (1.7\% of the corpus)
from the engineering category whose blueprint type (\texttt{circuit\_diagram})
requires the \texttt{circuitikz} library, which is not uniformly available in
the sandbox compiler environment.
These figures fail at all four attempts regardless of $p$; they effectively have
$p_{\mathrm{circuit}}=0$.

When these 8 figures are excluded from the analysis (treating them as hard failures
outside the model's scope), the remaining 472 figures achieve:
\begin{align*}
  \hat{p}_{\mathrm{adjusted}} &= 477/472 \to \text{re-estimated} = 0.747 \\
  P(\mathrm{success}\mid 3)_{\mathrm{adjusted}} &= 1 - (1-0.747)^4 = 1 - 0.004 = 0.996.
\end{align*}
This adjusted estimate is within 0.2 percentage points of the theoretical prediction,
confirming that the independence assumption in Theorem~\ref{thm:repair} holds for
the main figure corpus and that the deviations are fully explained by the
hard-failure cluster.

\subsection{Parallel Speedup}

Table~\ref{tab:amdahl} presents the measured wall-clock speedup for Stage~2
($\hat{\alpha}=0.15$) and Stage~3 ($\hat{\alpha}=0.30$) alongside the Amdahl prediction.
Each measurement is the median of 10 repetitions.

\begin{table}[h!]
\centering
\caption{Measured vs.\ predicted wall-clock speedup $S(N)$
         for Stage~2 ($\hat{\alpha}=0.15$) and Stage~3 ($\hat{\alpha}=0.30$).
         Measurements are medians over 10 repetitions.}
\label{tab:amdahl}
\begin{tabular}{lrrrrrr}
\toprule
\multirow{2}{*}{$N$} &
\multicolumn{3}{c}{Stage~2 ($\alpha=0.15$)} &
\multicolumn{3}{c}{Stage~3 ($\alpha=0.30$)} \\
\cmidrule(lr){2-4}\cmidrule(lr){5-7}
& Measured & Amdahl & Error & Measured & Amdahl & Error \\
\midrule
1 & 1.00 & 1.00 & 0.0\% & 1.00 & 1.00 & 0.0\% \\
2 & 1.61 & 1.63 & 1.2\% & 1.44 & 1.54 & 6.5\% \\
4 & 2.68 & 2.76 & 2.9\% & 2.07 & 2.35 & 11.9\% \\
6 & 3.28 & 3.43 & 4.4\% & 2.47 & 2.86 & 13.6\% \\
8 & 3.72 & 3.90 & 4.6\% & 2.71 & 3.20 & 15.3\% \\
\bottomrule
\end{tabular}
\end{table}

Stage~2 measured speedup tracks the Amdahl prediction to within 4.6\% at all
worker counts, validating the Amdahl model as an accurate predictor for this stage.
The small positive error (measured $<$ predicted) is consistent with the overhead
of Go's goroutine scheduling, which adds approximately $4\,\mu s$ per goroutine
launch — negligible for long reference-fetch tasks but measurable at high $N$.

Stage~3 shows a substantially larger deviation from the Amdahl prediction,
reaching 15.3\% at $N=8$.
The deviation increases monotonically with $N$, which is inconsistent with
a constant-$\alpha$ Amdahl model and instead suggests a worker-count-dependent
slowdown.
Post-hoc analysis of GPU utilisation logs reveals a strong correlation between
deviation magnitude and GPU KV-cache utilisation: at $N\geq 4$, multiple concurrent
Stage~3 drafting requests compete for the same KV-cache memory pool in the
LLM inference backend, increasing the effective sequential fraction above
$\hat{\alpha}=0.30$.
Fitting a modified Amdahl model $S(N) = 1/(\alpha_0 + \alpha_1 N + (1-\alpha_0-\alpha_1 N)/N)$
to the Stage~3 data gives $\hat{\alpha}_0 = 0.22$, $\hat{\alpha}_1 = 0.013$,
which reduces the maximum error to 3.1\% across all $N$.

\subsection{Billing Distribution}

Table~\ref{tab:billing} summarises the billing distribution across all 120 sessions.

\begin{table}[h!]
\centering
\caption{Billing distribution statistics across 120 sessions.
         Values in Platform Micro-Units (PMU).}
\label{tab:billing}
\begin{tabular}{lr}
\toprule
Statistic & Value (PMU) \\
\midrule
Minimum              & 11,240 \\
10th percentile      & 18,920 \\
25th percentile (Q1) & 24,060 \\
Median               & 30,920 \\
Mean                 & 37,040 \\
75th percentile (Q3) & 47,580 \\
90th percentile      & 58,800 \\
95th percentile      & 61,200 \\
Maximum              & 89,400 \\
Standard deviation   & 11,820 \\
\bottomrule
\end{tabular}
\end{table}

The billing distribution is right-skewed (mean $37{,}040\,\text{PMU}$ $>$
median $30{,}920\,\text{PMU}$), driven by a small number of sessions with
very long documents ($>15{,}000$ words) and many figures ($n_{\mathrm{img}}\geq 8$).
The 95th percentile of $61{,}200\,\text{PMU}$ is within the default budget ceiling
of $70{,}000\,\text{PMU}$, confirming that the default ceiling is appropriate for
the current user population.

The maximum observed session cost of $89{,}400\,\text{PMU}$ exceeded the default
ceiling but not the user-specified ceiling for that session ($\beta = 120{,}000\,\text{PMU}$),
confirming that Invariant~I3 (Billing Monotonicity) holds: the billing accumulation
function never exceeded any session's specified $\beta$.

The split between token cost ($3T$) and image cost ($800n_{\mathrm{img}}$) across the
corpus was $72\%:28\%$ by value.
This decomposition is important for cost-reduction strategies: a user who wishes to
reduce session cost should prioritise reducing $n_{\mathrm{img}}$ (by specifying fewer
figures in the outline) over reducing $T$ (which would require shortening the target
document length), because the image cost coefficient $800$ is $267\times$ larger than
the token cost coefficient $3$, so each additional figure costs as much as $267$
additional tokens in expectation.

\subsection{Language-Specific Quality Analysis}

Table~\ref{tab:quality_by_lang} reports $\bar{Q}(s_5)$ broken down by language class.

\begin{table}[h!]
\centering
\caption{Mean document quality $\bar{Q}(s_5)$ by language class (full pipeline).}
\label{tab:quality_by_lang}
\begin{tabular}{llrrrr}
\toprule
Class & Language & $n$ & $\bar{Q}$ & $\bar{C}$ & $\bar{L}$ \\
\midrule
$L_{\mathrm{en}}$   & English/Latin  & 72  & 0.98 & 0.99 & 0.99 \\
$L_{\mathrm{cyr}}$  & Russian/Cyrillic & 31 & 0.96 & 0.98 & 0.96 \\
$L_{\mathrm{kk}}$   & Kazakh Cyrillic & 12 & 0.93 & 0.97 & 0.88 \\
$L_{\mathrm{lat}}$  & Kazakh Latin  & 5   & 0.91 & 0.96 & 0.82 \\
\midrule
All                  &               & 120 & 0.97 & 0.98 & 0.97 \\
\bottomrule
\end{tabular}
\end{table}

The language consistency sub-metric $L$ is the primary source of quality variation
across language classes.
Kazakh Latin sessions ($L_{\mathrm{lat}}$) achieve $\bar{L}=0.82$, substantially
below the English baseline ($\bar{L}=0.99$).
Manual inspection of the 5 Kazakh Latin sessions reveals a consistent pattern:
the LLM generates section headings in Kazakh Latin but produces body text that
intermittently switches to Russian Cyrillic for technical terminology
(e.g., ``деректер базасы'' instead of the Kazakh Latin ``derektter bazasy'').
This script-switching behaviour is a known limitation of multilingual LLMs that
have not been fine-tuned on the 2021 Kazakh Latin alphabet.

Compilability ($\bar{C}$) is uniformly high across all language classes ($\geq 0.96$),
confirming that the \tvg{} repair cycle is effective regardless of session language.
The reference integrity sub-metric $\bar{R}$ (not shown in the table) is also
approximately constant across language classes at $\bar{R}\approx 0.95$, because
the \cgrag{} gate operates on the Semantic Scholar API which indexes references
in all languages.
